% ============================================================
%  MockAI – Full 44-Page Project Report (LaTeX)
%  Department of AI & ML, Kongu Engineering College
%  October 2026
% ============================================================
\documentclass[12pt,a4paper]{report}

% ---- Packages -----------------------------------------------
\usepackage[top=2.5cm,bottom=2.5cm,left=3.5cm,right=2.5cm]{geometry}
\usepackage{times}
\usepackage{setspace}
\usepackage{graphicx}
\usepackage{amsmath,amssymb}
\usepackage{array}
\usepackage{booktabs}
\usepackage{longtable}
\usepackage{hyperref}
\usepackage{listings}
\usepackage{xcolor}
\usepackage{caption}
\usepackage{fancyhdr}
\usepackage{titlesec}
\usepackage{tocloft}
\usepackage{enumitem}
\usepackage{parskip}
\usepackage{indentfirst}
\usepackage{microtype}

% ---- Spacing & Formatting ------------------------------------
\onehalfspacing
\setlength{\parindent}{1.25cm}

% ---- Chapter Title Style ------------------------------------
\titleformat{\chapter}[block]
  {\centering\bfseries\large}
  {CHAPTER \thechapter}
  {0pt}
  {\\\vspace{6pt}\bfseries\large\MakeUppercase}

\titleformat{\section}[block]
  {\bfseries\normalsize}
  {\thesection}
  {1em}{}

\titleformat{\subsection}[block]
  {\bfseries\normalsize}
  {\thesubsection}
  {1em}{}

% ---- Header / Footer ----------------------------------------
\pagestyle{fancy}
\fancyhf{}
\fancyhead[R]{\thepage}
\fancyhead[L]{\slshape\leftmark}
\renewcommand{\headrulewidth}{0.4pt}

% ---- Code Listing Style -------------------------------------
\lstset{
  basicstyle=\ttfamily\footnotesize,
  breaklines=true,
  frame=single,
  numbers=left,
  numberstyle=\tiny,
  keywordstyle=\color{blue},
  commentstyle=\color{gray},
  stringstyle=\color{orange},
  showstringspaces=false
}

% ============================================================
\begin{document}

% ============================================================
%  PAGE 1 – TITLE PAGE
% ============================================================
\begin{titlepage}
\centering
\vspace*{1cm}
{\Large\bfseries MOCKAI: A MULTI-AGENT LLM FRAMEWORK FOR\\[6pt]
AUTOMATED TECHNICAL INTERVIEW ASSESSMENT WITH\\[6pt]
REAL TIME GAZE-BASED PROCTORING\\[1.5cm]}

{\large A PROJECT REPORT}\\[0.4cm]
{\large Submitted by}\\[1cm]

{\large\bfseries RITHIK R}\\
{\large 23ALR077}\\[0.8cm]

{\large\bfseries SANJAY K}\\
{\large 23ALR086}\\[1.5cm]

{\normalsize in partial fulfilment of the requirements\\
for the award of the degree of}\\[0.4cm]
{\large\bfseries BACHELOR OF TECHNOLOGY}\\
{\large IN}\\[0.4cm]
{\large\bfseries ARTIFICIAL INTELLIGENCE AND MACHINE LEARNING}\\[0.8cm]

{\large DEPARTMENT OF ARTIFICIAL INTELLIGENCE AND\\
MACHINE LEARNING}\\[0.4cm]

{\normalsize (Autonomous)}\\[1cm]

\includegraphics[width=3cm]{logo.png}\\[0.8cm]

{\large PERUNDURAI ERODE -- 638~060}\\[0.4cm]
{\large OCTOBER 2026}
\end{titlepage}

% ============================================================
%  PAGE 2 – BONAFIDE CERTIFICATE
% ============================================================
\clearpage
\thispagestyle{plain}
\begin{center}
{\large\bfseries DEPARTMENT OF ARTIFICIAL INTELLIGENCE AND\\
MACHINE LEARNING\\
KONGU ENGINEERING COLLEGE\\
(Autonomous)\\
PERUNDURAI ERODE -- 638060\\
OCTOBER 2026}
\end{center}

\vspace{1.5cm}
\begin{center}
{\large\bfseries BONAFIDE CERTIFICATE}
\end{center}

\vspace{0.8cm}
\noindent This is to certify that the project report entitled \textbf{MOCKAI: A MULTI-AGENT LLM
FRAMEWORK FOR AUTOMATED TECHNICAL INTERVIEW ASSESSMENT
WITH REAL-TIME GAZE-BASED PROCTORING} is the Bonafide record of project
work done by \textbf{RITHIK R\,(23ALR077)}, \textbf{SANJAY K\,(23ALR086)} in partial
fulfilment of the requirements for the award of the Degree of Bachelor of Engineering in
Anna University Chennai during the year 2025--2026.

\vspace{2.5cm}
\noindent\begin{tabular}{@{}p{7cm}@{}p{7cm}@{}}
\textbf{SUPERVISOR} & \textbf{HEAD OF THE DEPARTMENT}\\[0.3cm]
(Signature with seal) & (Signature with seal)\\[2cm]
Date:\underline{\hspace{3cm}} &
\end{tabular}

\vspace{1.5cm}
\noindent Submitted for the end semester viva voce examination held on
\underline{\hspace{5cm}}

\vspace{2.5cm}
\noindent\begin{tabular}{@{}p{7cm}@{}p{7cm}@{}}
\textbf{INTERNAL EXAMINER} & \textbf{EXTERNAL EXAMINER}
\end{tabular}
\addcontentsline{toc}{chapter}{BONAFIDE CERTIFICATE}

% ============================================================
%  PAGE 3 – DECLARATION
% ============================================================
\clearpage
\thispagestyle{plain}
\begin{center}
{\large\bfseries DEPARTMENT OF ARTIFICIAL INTELLIGENCE AND\\
MACHINE LEARNING\\
KONGU ENGINEERING COLLEGE\\
(Autonomous)\\
PERUNDURAI ERODE -- 638060\\
OCTOBER 2026}
\end{center}

\vspace{1.2cm}
\begin{center}
{\large\bfseries DECLARATION}
\end{center}

\vspace{0.8cm}
\noindent We affirm that the Project Report titled \textbf{MOCKAI: A MULTI-AGENT LLM
FRAMEWORK FOR AUTOMATED TECHNICAL INTERVIEW ASSESSMENT
WITH REAL-TIME GAZE-BASED PROCTORING} being submitted in partial
fulfilment of the requirements for the award of Bachelor of Engineering is the original
work carried out by us. It has not formed the part of any other project report or
dissertation on the basis of which a degree or award was conferred on an earlier occasion
on this or any other candidate.

\vspace{2.5cm}
\noindent Date:\underline{\hspace{3cm}}

\vspace{1cm}
\noindent\begin{tabular}{@{}p{7cm}@{}p{6cm}@{}}
& \textbf{RITHIK R}\\
& (23ALR077)\\[1.2cm]
& \textbf{SANJAY K}\\
& (23ALR086)
\end{tabular}

\vspace{1.2cm}
\noindent I certify that the declaration made by the above candidates is true to the best
of my knowledge.

\vspace{2cm}
\noindent Date:\underline{\hspace{3cm}}

\vspace{0.5cm}
\noindent Name and Signature of the Supervisor with seal
\addcontentsline{toc}{chapter}{DECLARATION}

% ============================================================
%  PAGE 4 – ACKNOWLEDGEMENT
% ============================================================
\clearpage
\thispagestyle{plain}
\begin{center}
{\large\bfseries ACKNOWLEDGEMENT}
\end{center}
\addcontentsline{toc}{chapter}{ACKNOWLEDGEMENT}

\vspace{0.5cm}
\noindent We express our sincere thanks and gratitude to Thiru.\ E.\ R.\ K.\ KRISHNAN, M.Com.,
our beloved Correspondent and all other philanthropic trust members of the Kongu Vellalar
Institute of Technology Trust who have always encouraged us in the academic and
co-curricular activities.

\vspace{0.4cm}
\noindent We are extremely thankful to the dynamic Principal Dr.\ R.\ PARAMESHWARAN, M.E.,
Ph.D., for providing the necessary facilities to complete our work.

\vspace{0.4cm}
\noindent We would like to express our sincere gratitude to our respected Head of the Department
Dr.\ T.\ ABIRAMI, MCA., M.E., Ph.D., for providing the necessary facilities.

\vspace{0.4cm}
\noindent We extend our thanks to N.\ KANIMOZHI, M.E., (Ph.D)., Assistant Professor, Department
of Artificial Intelligence and Machine Learning, Project Coordinator for her encouragement
and valuable advice that made us carry out the project work successfully.

\vspace{0.4cm}
\noindent We extend our gratitude to N.\ RENUKA, M.Tech., Assistant Professor, Department of
Artificial Intelligence and Machine Learning, who supervised with valuable ideas and
suggestions, which have been very helpful in the project. We are grateful to all the faculty
members of the Artificial Intelligence and Machine Learning Department for their support.

\vspace{0.4cm}
\noindent We are grateful to KONGU ENGINEERING COLLEGE for giving us an opportunity to
learn new skills and technologies and enhance our knowledge by providing the necessary
infrastructure.

% ============================================================
%  PAGE 5 – ABSTRACT
% ============================================================
\clearpage
\thispagestyle{plain}
\begin{center}
{\large\bfseries ABSTRACT}
\end{center}
\addcontentsline{toc}{chapter}{ABSTRACT}

\vspace{0.5cm}
\noindent Technical recruitment and campus placement systems require intelligent, automated
assessment platforms that provide scalable, objective, and multi-dimensional candidate
evaluation. Although Large Language Models (LLMs) demonstrate strong conversational
reasoning, relying on single models in isolation results in unmonitored evaluation
environments, unsafe code execution, and a lack of anti-hallucination guardrails. To address
these challenges, this project proposes \textbf{MockAI}, an autonomous AI-powered technical
assessment and real-time proctoring platform that integrates Retrieval-Augmented Generation
(RAG), agentic workflows via LangGraph, edge computer vision supervision, and
Docker-isolated code execution.

\vspace{0.4cm}
\noindent The proposed system employs a multi-round evaluation pipeline comprising general
aptitude, dynamic domain-specific technical interviews, and automated Applicant Tracking
System (ATS) resume screening. Real-time proctoring is achieved through an edge computer
vision engine utilizing OpenCV Haar Cascades and spatial moment analysis for gaze
direction estimation ($0.35 < \text{RelativeX} < 0.65$), tracking face presence and
prolonged blinks ($>$1.5\,s) with automated multi-strike warning governance. For
multi-source technical assessments, a four-node LangGraph pipeline ingests PDFs,
YouTube video transcripts, and audio files, vectorizing chunks into FAISS and ChromaDB
using Sentence Transformers (\texttt{all-MiniLM-L6-v2}) to generate targeted assessments
via MultiQuery retrieval. Candidate coding capability is evaluated through an air-gapped
Docker sandbox (NetworkMode: `none', 128\,MB RAM, 50\% CPU quota) supporting C, C++,
Java, and Python.

\vspace{0.4cm}
\noindent We built MockAI using a decoupled microservices setup, pairing a React~19 and Monaco
Editor interface with a Flask backend, Node.js container management, Groq LLaMA-3.1,
and MongoDB Atlas. To keep technical responses accurate, a four-stage LangGraph agent
filters off-topic prompts and restructures queries before retrieving data from Pinecone.
Across intensive interview simulations and code benchmarks, the system proved both fast
and secure---delivering proctoring latency under 180\,ms, maintaining 100\% sandbox
containment against malicious code, and achieving a 92.4\% match accuracy for resume
screening. These results demonstrate that combining AI agent workflows with isolated
sandboxes and vision proctoring creates a secure, reliable hiring platform, paving the way
for future updates like voice analysis and adaptive testing.

% ============================================================
%  TABLE OF CONTENTS, LIST OF TABLES, LIST OF FIGURES
% ============================================================
\clearpage
\renewcommand{\contentsname}{TABLE OF CONTENTS}
\tableofcontents

\clearpage
\renewcommand{\listtablename}{LIST OF TABLES}
\listoftables
\addcontentsline{toc}{chapter}{LIST OF TABLES}

\clearpage
\renewcommand{\listfigurename}{LIST OF FIGURES}
\listoffigures
\addcontentsline{toc}{chapter}{LIST OF FIGURES}

% ============================================================
%  PAGE 10 – CHAPTER 1: INTRODUCTION
% ============================================================
\chapter{INTRODUCTION}
\pagenumbering{arabic}
\setcounter{page}{1}

MockAI is an advanced Artificial Intelligence-driven technical assessment and interview
platform designed to improve the objectivity, scalability, and integrity of the recruitment
process. The system combines conversational Large Language Models (LLMs), agentic
workflows via LangGraph, real-time computer vision proctoring, and Docker-isolated code
execution to provide candidates and evaluators with a holistic assessment environment
instead of relying solely on manual, subjective interviewers or static multiple-choice tests.
By orchestrating multi-round interview pipelines and continuous visual monitoring before
compiling performance insights, MockAI reduces human bias and assessment dishonesty
while providing candidates with structured, quantifiable feedback grounded in candidate
performance.

MockAI replaces static, unmonitored testing with a multi-tiered, intelligent evaluation
pipeline. It orchestrates a three-round interview workflow covering general aptitude,
technical inquiry, and behavioral HR assessments. Powered by Groq-accelerated LLaMA-3.1
models, the conversational engine maintains session memory throughout the evaluation. For
practical programming, an air-gapped Docker sandbox safely compiles and executes
candidate code across multiple languages. A LangGraph state-graph pipeline dynamically
ingests multi-source learning materials into FAISS and ChromaDB vector stores to generate
targeted quizzes. The complete system is implemented using React for the frontend
interface, Express.js for request routing, Flask for backend orchestration, ChromaDB for
semantic retrieval, and MongoDB Atlas for maintaining candidate history.

\section{EXISTING SYSTEM}

Currently, most technical interview and screening portals rely on static questionnaires,
pre-recorded videos, or generic question banks. While capable of basic testing, they fail to
deliver realistic and objective candidate evaluations. Lacking conversational reasoning and
adaptive probing, these systems cannot evaluate complex decisions.

Additionally, manual interviews suffer from human fatigue and unconscious bias, whereas
automated platforms offer non-interactive assessments that lack actionable performance
feedback.

Existing recruitment tools also lack continuous real-time proctoring and isolated execution
environments for code evaluation. Consequently, unmonitored webcam assessments remain
vulnerable to academic dishonesty, unauthorized external screen usage, and candidate
substitution. Without containerized sandboxing, executing candidate scripts exposes host
infrastructure to resource exhaustion, infinite loops, and security vulnerabilities, while
traditional Applicant Tracking Systems (ATS) rely on rigid keyword matching that penalizes
qualified applicants. These limitations highlight the need for a secure, intelligent, and
multi-dimensional assessment framework.

Our proposed MockAI system addresses these shortcomings by unifying multi-round
conversational LLM evaluation with session memory, edge computer vision proctoring,
Docker-isolated multi-language code execution, agentic LangGraph workflows, and
semantic ATS resume parsing. This combination ensures an objective, secure, and traceable
candidate assessment pipeline featuring dynamic, domain-specific questioning and
continuous real-time proctoring throughout the recruitment lifecycle.

\section{SYSTEM STUDY}

The system study examines the practical requirements, existing limitations, user
expectations, and functional scenarios necessary for developing an intelligent technical
interview assessment system. Understanding these factors helps ensure that MockAI
addresses real-world recruitment challenges while providing reliable and explainable
assistance to candidates and recruiters.

\subsection{Business Requirements}

MockAI is designed as an autonomous, multi-tiered interview and technical screening
platform capable of supporting evidence-based candidate evaluations across theoretical,
practical, and behavioral domains. The system must dynamically generate domain-specific
interview questions, enforce continuous video-based proctoring, execute untrusted user
code within isolated container environments, parse resumes against target job descriptions,
and extract quiz content from diverse multimedia sources. The primary users include
students, job seekers, academic placement coordinators, and enterprise HR recruiters. The
system should provide an intuitive interface, maintain secure evaluation history, and
generate structured numerical reports with actionable feedback while remaining scalable
for future recruitment applications.

\subsection{Analysis of Existing System}

The analysis of existing interview preparation platforms reveals several limitations
affecting transparency and reliability. Most systems generate static questions directly from
non-adaptive banks or unmonitored online forms without real-time gaze tracking or secure
execution environments. As a result, unmonitored tests are vulnerable to academic
dishonesty, while code assessments risk host infrastructure corruption due to a lack of
container isolation. Users generally cannot receive detailed behavioral scoring, automated
skill-gap analysis, or real-time proctoring feedback.

MockAI improves upon these limitations by introducing conversational LLM memory
through Groq, edge computer vision proctoring via OpenCV, Docker-isolated multi-language
sandboxing, agentic LangGraph workflows, and MongoDB storage. These improvements
create a more objective evaluation process while maintaining security throughout the
recruitment pipeline.

\subsection{Identifying User Needs}

Users expect technical assessment systems to provide interview evaluations that are fair
and realistic while remaining secure and informative. They require actionable performance
feedback rather than simple rejection notices, clear proctoring alerts that preserve test
integrity, and the ability to test code safely in real time. Candidates and recruiters also
benefit from automated resume matching and multi-source quiz generation instead of
static question lists.

MockAI addresses these requirements by combining conversational LLM orchestration,
real-time gaze tracking, containerized execution, and semantic resume parsing. The system
also supports persistent result tracking, allowing candidate growth and historical
performance to be monitored across multiple interview rounds.

\subsection{Developing Use Cases}

MockAI supports several practical recruitment-oriented scenarios that demonstrate its
usefulness across educational and enterprise environments. A candidate may initiate a mock
interview session, after which the system verifies user identity, poses domain-specific
technical questions, records responses with conversation memory, tracks pupil gaze, and
evaluates code inside an isolated Docker sandbox. If gaze deviations occur, the proctoring
engine issues real-time warnings before compiling the final score report.

Additional use cases include aptitude testing, resume screening, quiz generation from study
materials, and candidate performance analytics. These scenarios demonstrate the flexibility
of MockAI in providing transparent technical reasoning while maintaining an explainable
decision-making process suitable for future intelligent recruitment applications.

\section{OBJECTIVE}

The objective of this project is to design and implement MockAI, an autonomous technical
assessment and interview platform that combines conversational Large Language Models
with real-time vision proctoring and sandboxed code execution to improve scalability,
security, and objectivity in candidate evaluation systems. The proposed framework aims
to replace subjective manual interviews and unmonitored online testing with an end-to-end,
transparent recruitment evaluation suite.

The system employs a three-round interview pipeline covering general aptitude, technical
inquiry, and behavioral HR rounds using Groq-accelerated LLaMA-3.1 models equipped
with session memory. An edge computer vision engine enforces real-time gaze proctoring,
an air-gapped Docker sandbox compiles multi-language candidate code under resource
limits, and a four-stage LangGraph pipeline ingests multi-source educational materials into
FAISS and ChromaDB vector stores.

The project is evaluated using simulated interview sessions, container security stress tests,
and computer vision gaze-tracking benchmarks, allowing the effectiveness of agentic LLM
orchestration and sandboxed evaluation to be examined through measurable experimental
performance.

\section{SCOPE}

The scope of this project encompasses the complete development and evaluation of an
autonomous technical recruitment and assessment system integrating document processing,
vector retrieval, conversational interview workflows, edge vision proctoring, sandboxed
code execution, and ATS resume matching.

The system begins by ingesting multi-source educational materials and resumes,
normalizing textual content, generating vector embeddings with Sentence Transformers
(\texttt{all-MiniLM-L6-v2}), and indexing them within Pinecone, FAISS, and ChromaDB
for semantic retrieval during assessment generation.

Architecturally, MockAI combines React~19 and Tailwind CSS for the user interface,
Monaco Code Editor for code editing, Flask~3.0 for backend API routing, Node.js with
Dockerode for container management, and MongoDB Atlas for persistence. The proctoring
pipeline employs OpenCV Haar Cascades and spatial moments to track pupil gaze direction
with automated strike governance, while the Docker sandbox executes C, C++, Java, and
Python code under complete network isolation and memory constraints. The LangGraph
state-graph pipeline dynamically restructures queries and filters irrelevant requests before
synthesizing adaptive quizzes.

The project also includes experimental evaluation across proctoring latency, container
isolation containment, and ATS match accuracy, establishing persistent candidate records
for administrative placement dashboards. As an AI-driven recruitment prototype, MockAI
establishes a scalable foundation for future intelligent assessment applications while
emphasizing security, fairness, and objective candidate evaluation.

% ============================================================
%  PAGE 15 – CHAPTER 2: LITERATURE REVIEW
% ============================================================
\chapter{LITERATURE REVIEW}

The development of automated technical assessment and mock interview platforms has
rapidly advanced through the integration of Large Language Models (LLMs),
Retrieval-Augmented Generation (RAG), containerized code virtualization, and real-time
computer vision proctoring. Recent studies have focused on improving the objectivity and
engagement of candidate evaluations by replacing static question forms with adaptive,
stateful dialogue agents, safeguarding exam integrity via non-intrusive edge gaze tracking,
and grounding assessment material in verified technical curricula. This chapter reviews the
major research works that form the theoretical foundation of MockAI and explains how
these contributions directly influenced the proposed architecture.

\section{REVIEW OF EXISTING RESEARCH}

In~[1], Brown et al.\ demonstrated that scaling autoregressive language models enables
few-shot and zero-shot reasoning across diverse natural language tasks without task-specific
fine-tuning. The study proved that general-purpose language models can comprehend
complex prompts, evaluate context, and generate coherent dialogue, serving as the core
linguistic foundation for the dynamic multi-round interview agents deployed in MockAI.

In~[2], Nigam et al.\ presented a systematic survey of automated online proctoring systems,
categorizing methodologies across facial verification, gaze estimation, hardware peripheral
detection, and acoustic monitoring. The authors identified critical limitations in existing
systems, including high computational overhead, intrusive installation requirements, and
elevated false-positive rates from rigid gaze boundaries, directly motivating MockAI's
lightweight, edge computer vision pipeline.

In~[3], Pal et al.\ investigated context-aware technical question generation, demonstrating
that condition-specific language models synthesize relevant and non-repetitive inquiries
when primed with candidate metadata. Their findings inspired MockAI's dynamic
questioning mechanism that tailors interview questions according to a candidate's
self-declared domain competencies and project history.

In~[4], Rietz and Maedche developed Ladderbot, a conversational agent for semi-structured
online laddering interviews. The study showed that agents maintaining stateful dialogue
history and adaptive follow-up probing achieve higher candidate engagement and response
depth than static web surveys, validating MockAI's three-round interview workflow across
introductory, project deep-dive, and core subject technical rounds.

In~[5], Lewis et al.\ introduced Retrieval-Augmented Generation (RAG), combining
parametric language models with non-parametric dense vector index retrieval. By retrieving
authoritative document passages prior to response generation, RAG mitigates factual
hallucinations, establishing the foundational strategy used in MockAI's multi-source
LangGraph quiz generator and portfolio evaluation modules.

In~[6], Atoum et al.\ proposed a multi-component online exam monitoring system using
visual sensors to detect unauthorized behavior such as eye divergence, head rotation, and
candidate absence. Their experimental results proved that temporal monitoring across
continuous intervals provides robust academic integrity verification, which directly
informed MockAI's temporal strike-governed proctoring state machine.

In~[7], Murthy and Biswas explored appearance-based gaze estimation techniques utilizing
spatial difference mechanisms to track ocular gaze from monocular webcam feeds. Their
findings showed that calculating pupil displacement relative to the eye bounding box
yields dependable visual focus vectors without requiring expensive hardware, directly
supporting MockAI's pupil centroid extraction using spatial image moments
($M_{10} / M_{00}$).

In~[8], Viola and Jones introduced the classic object detection framework utilizing
Haar-like features, integral images, and AdaBoost classifier cascades. Their methodology
enables rapid facial and ocular localization with minimal computational footprint on
standard hardware, forming the primary face and eye detection backbone within MockAI's
real-time edge proctoring pipeline.

In~[9], Topsakal and Akinci presented a systematic framework for building language model
applications using LangChain abstractions. The study detailed how prompt templates,
chains, and conversation buffer memories maintain contextual continuity across extended
interactions, guiding the implementation of MockAI's stateful interview memory store.

In~[10], Wu et al.\ introduced AutoGen, a multi-agent framework enabling autonomous
language model agents to communicate, delegate sub-tasks, and execute complex workflows
cooperatively. Their findings highlighted that decomposing complex tasks into specialized,
role-governed agents improves factual precision, establishing the architectural basis for
MockAI's specialized interview and evaluation agents.

In~[11], Liang et al.\ explored multi-agent critique frameworks to foster factual convergence
in language models. The authors showed that subjecting candidate responses to multi-agent
evaluation and rubric consensus mitigates scoring variance, inspiring MockAI's
rubric-driven evaluation chain that independently evaluates strengths, weaknesses, and
numerical marks out of 50.

In~[12], Zhang et al.\ developed Multi-Task Cascaded Convolutional Networks (MTCNN)
for joint face detection and landmark alignment, demonstrating the importance of
normalized facial boundary extraction prior to downstream biometric tracking. Their work
emphasized precise facial ROI isolation before assessing eye regions, directly informing
MockAI's two-stage face-to-eye localization flow.

\section{LITERATURE SUMMARY}

The reviewed literature demonstrates that recent advancements in Artificial Intelligence
have significantly improved the capabilities of Large Language Models, Retrieval-Augmented
Generation, computer vision proctoring, multi-agent coordination, and sandboxed code
execution. State-of-the-art conversational models such as LLaMA-3.1 can conduct nuanced
technical dialogues~[1],~[4], while RAG pipelines effectively ground assessment generation
in syllabus documents and course materials~[5],~[9]. Concurrently, real-time computer
vision algorithms operating on facial landmarks and ocular centroids offer non-intrusive
mechanisms to verify candidate visual focus during remote assessments~[6],~[7],~[8],~[12].

Despite these individual advancements, existing technical assessment platforms suffer from
significant architectural fragmentation. Online testing portals typically deploy isolated
environments where practical coding workspaces lack continuous visual proctoring, or
remote surveillance tools operate merely as passive camera loggers without interactive
feedback~[2].

Furthermore, conventional automated interviewers rely on single-prompt evaluations that
lack persistent conversation memory, fail to sandbox untrusted candidate code execution
safely, and do not provide actionable feedback on technical skill deficiencies~[3],~[10].

MockAI bridges these critical gaps by synthesizing these disparate domains into a unified,
multi-agent assessment framework. By coupling LangChain stateful conversation
memory~[9] and multi-agent rubric evaluation~[10],~[11] with real-time Haar Cascade
spatial moment pupil tracking~[7],~[8], air-gapped Docker container isolation, ATS resume
analysis, and LangGraph-orchestrated multi-source quiz generation~[5], MockAI delivers
an objective, secure, and fully automated recruitment and interview preparation platform.

% ============================================================
%  PAGE 19 – CHAPTER 3: REQUIREMENTS
% ============================================================
\chapter{REQUIREMENTS}

The successful implementation of MockAI requires suitable hardware resources, software
tools, and development frameworks that work together to support real-time video
processing, containerized code execution, agentic LLM orchestration, and web-based
interaction. The system is designed as a modern full-stack application where the frontend,
backend gateway, container execution microservice, vector databases, and language model
services communicate efficiently to deliver an objective and tamper-resistant assessment
ecosystem. This chapter describes the hardware and software requirements needed for
development and explains the major technologies used in building the proposed system.

\section{HARDWARE REQUIREMENTS}

MockAI is developed as a web-based application that performs edge computer vision
proctoring, multi-round interview dialogues, multi-source document ingestion, and isolated
code execution. Since the system uses cloud-based language model providers such as Groq
Cloud API for high-throughput LLaMA-3.1 inference, high-end local GPU hardware is not
essential for normal operation. However, a standard development computer with sufficient
memory and multi-core processing capability is required to run the frontend, backend
gateway, container engine, and local storage simultaneously.

A development system featuring an Intel Core i5 or AMD Ryzen~5 processor (4~cores,
8~threads) ensures smooth overall performance. At least 8\,GB of RAM (16\,GB
recommended) and 15\,GB of SSD storage support concurrent multi-container execution.
A standard 720p 30\,FPS webcam is required to capture video streams for real-time
computer vision proctoring. A stable internet connection enables reliable communication
with external Groq LLM endpoints and Pinecone vector store.

\section{SOFTWARE REQUIREMENTS}

MockAI combines several software technologies to create an automated technical assessment
and interview preparation platform. Python~3.9+ serves as the primary programming
language for backend development, while Flask~3.0 handles lightweight REST API routing,
request validation, and modular blueprint management. Node.js with Express~5.1 provides
the container management microservice that interfaces with the Docker daemon via
Dockerode. React~19 with Vite~6 forms the client-side single-page application, while
MongoDB Atlas provides cloud-hosted persistent storage. The LangGraph library powers
the agentic state-graph pipeline, and Groq Cloud supplies accelerated LLaMA-3.1
inference without requiring local GPU resources.

\section{SOFTWARE DESCRIPTION}

MockAI integrates multiple software components that work together to perform visual
proctoring, conversational interview dialogues, sandboxed code execution, and knowledge
retrieval. Flask~3.0 acts as the central REST API gateway that coordinates system
communication, validates requests, and routes traffic across dedicated blueprints including
\texttt{interview}, \texttt{facetrack}, \texttt{quiz}, \texttt{ats}, \texttt{questionbank},
and \texttt{Domainforum}. React~19 provides a dynamic interface where candidates navigate
assessments, interact with the Monaco Code Editor (\texttt{@monaco-editor/react}), observe
proctoring indicators, and view performance evaluations.

Express.js and Node.js function as the container management microservice that connects
with the host Docker daemon through Dockerode to provision ephemeral, air-gapped
sandboxes with strict resource limits for safe code compilation. Educational materials and
documents are processed using PyPDF2, python-docx, and transcript APIs, after which
Sentence Transformers generate 384-dimensional semantic embeddings stored in FAISS,
ChromaDB, and Pinecone for Top-$K$ retrieval during quiz generation. MongoDB Atlas
maintains a persistent record of candidate profiles, job roles, and aggregated multi-round
test marks to support transparency and auditability throughout every session.

The intelligence layer is powered by Groq Cloud LLaMA-3.1 models using LangChain
\texttt{ConversationBufferMemory} to maintain session context across Aptitude, Technical,
and HR interview rounds, yielding a score out of 50. OpenCV processes base64 webcam
frames using contour spatial moments ($M_{10}/M_{00}$) to compute pupil displacement
($0.35 < \text{RelativeX} < 0.65$) and enforce an automated three-strike proctoring
warning protocol. Finally, a LangGraph state-graph pipeline cleans filler words from
ingested learning sources and employs a \texttt{MultiQueryRetriever} to synthesize
targeted multiple-choice question assessments.

\section{USER INTERFACE}

The MockAI user interface is designed to provide a simple, responsive, and transparent
experience that allows candidates and HR administrators to observe every stage of the
technical evaluation process. Unlike conventional testing portals that display static
question lists, the interface presents real-time proctoring feedback, an interactive code
editor, dynamic interview prompts, and automated resume analysis in a structured
dashboard.

Candidates begin by authenticating and selecting a job role through the setup interface,
after which the system verifies webcam alignment and initiates the three-round interview
workflow. The conversational AI interface presents dynamic questions alongside real-time
proctoring status alerts, while the technical coding workspace features a split-screen
layout with problem constraints and the Monaco Code Editor. For resume analysis, the ATS
module renders immediate match percentages, skill-gap badges, and structural
recommendations based on candidate uploads.

The dashboard-based design improves usability by organizing complex AI and proctoring
workflows into an understandable user journey, enabling candidates and recruitment teams
to navigate the complete process from interview initiation to quantitative report generation
while maintaining transparency and ease of operation.

% ============================================================
%  PAGE 22 – CHAPTER 4: PROPOSED SYSTEM
% ============================================================
\chapter{PROPOSED SYSTEM}

MockAI is an autonomous, multi-agent artificial intelligence framework and technical
assessment platform that integrates conversational Large Language Models (LLMs),
real-time edge computer vision proctoring, sandboxed multi-language code execution, and
agentic Retrieval-Augmented Generation (RAG). Unlike traditional assessment portals that
depend on unmonitored multiple-choice questionnaires or static question banks, the
proposed system dynamically conducts multi-round adaptive interviews, validates untrusted
candidate code within air-gapped Docker containers, and enforces continuous visual
integrity through gaze-tracking and blink monitoring. Every assessment metric, conversation
turn, and violation event is logged transparently, enabling candidates to receive granular
feedback while granting recruiters objective candidate evaluations. The system follows a
modular microservices architecture where client-side interaction, REST routing, video
analysis, code virtualization, and language model orchestration operate collaboratively to
deliver an end-to-end recruitment environment.

\section{SYSTEM ARCHITECTURE}

The MockAI architecture consists of interconnected modules that process candidate video
telemetry, natural language responses, uploaded resumes, and programming submissions
through dedicated pipelines before synthesizing the final performance evaluation. Initially,
the candidate accesses the platform through the React~19 single-page interface styled with
Tailwind CSS. The backend receives requests through Flask~3.0 and authenticates candidate
or HR credentials using salted bcrypt hashing. When an assessment begins, the client
webcam captures video frames that are continuously evaluated by an OpenCV computer
vision service to monitor facial presence, pupil gaze displacement, and blink intervals in
real time.

For conversational evaluations, the system coordinates an adaptive three-round interview
pipeline (Aptitude, Technical, and HR rounds) powered by Groq-accelerated LLaMA-3.1
inference. The engine maintains dialogue context via LangChain
\texttt{ConversationBufferMemory}, allowing follow-up technical questions to be
synthesized dynamically based on the candidate's declared project experiences.

\begin{figure}[htbp]
  \centering
  \includegraphics[width=\textwidth]{figures/fig1.png}
  \caption{MockAI System Architecture}
  \label{fig:system-arch}
\end{figure}

\section{MODULE DESCRIPTION}

The MockAI framework is organized into ten functional modules, each executing a
specialized responsibility within the recruitment, evaluation, and proctoring lifecycle.
These modules work together to ensure that every evaluation is grounded in objective
evidence, tamper-resistant monitoring, and comprehensive analytical scoring.

\subsection{CANDIDATE AUTHENTICATION \& ROLE MANAGEMENT}

The authentication module manages two user roles: HR Faculty (who define job roles,
register candidates, and view scorecards) and Students (who register, complete the
three interview rounds, and receive scored reports). Passwords are bcrypt-hashed with
12 salt rounds. On login, a session token is stored in \texttt{localStorage} and the React
\texttt{AuthContext}. Role-protected Higher-Order Components enforce route-level access
control, and all MongoDB queries are scoped to the logged-in HR Faculty member's email
to prevent cross-account data access.

\begin{table}[htbp]
  \centering
  \caption{Evaluation Domains \& Assessment Configurations}
  \label{tab:domains}
  \renewcommand{\arraystretch}{1.3}
  \begin{tabular}{p{4cm} p{9cm}}
    \toprule
    \textbf{Interview Domain} & \textbf{Assessment Topics Covered} \\
    \midrule
    Machine Learning \& AI (12) &
      Supervised learning, neural networks, CNNs, RNNs, transformers,
      LLMs, RAG, fine-tuning, overfitting, regularization, evaluation
      metrics, model deployment \\
    \addlinespace
    Backend Development (8) &
      REST APIs, Flask routing, authentication, MongoDB queries,
      bcrypt hashing, CORS, async handlers, microservices \\
    \addlinespace
    Data Structures \& Algorithms (10) &
      Arrays, linked lists, trees, graphs, dynamic programming,
      sorting, searching, time complexity, space complexity, recursion \\
    \addlinespace
    Computer Science Fundamentals (6) &
      Operating systems, deadlocks, caching, memory management,
      process scheduling, concurrency \\
    \addlinespace
    Docker \& DevOps (4) &
      Container isolation, NetworkMode, resource limits, CI/CD \\
    \addlinespace
    Behavioral / HR Round (4) &
      Conflict management, leadership, adaptability, teamwork \\
    \bottomrule
  \end{tabular}
\end{table}

\subsection{VIDEO INGESTION \& FACE ALIGNMENT}

The video ingestion module captures webcam frames from the candidate's client device at
standard video resolutions (720p at 30\,FPS). Frames are base64-encoded and transmitted
to the Flask proctoring endpoint (\texttt{/facetrack/process-frame}). The OpenCV engine
converts incoming images into grayscale representations and applies the pre-trained
\texttt{haarcascade\_frontalface\_default.xml} classifier to isolate candidate facial
bounding boxes $(x, y, w, h)$. Once a face region is localized, the pipeline extracts the
sub-region of interest (ROI) corresponding to the ocular region and applies
\texttt{haarcascade\_eye.xml} to detect left and right eyes, normalizing dimensions for
pupil analysis.

\subsection{PUPIL GAZE TRACKING \& BLINK ANALYSIS}

Once ocular bounding boxes are isolated, the gaze tracking module calculates the relative
horizontal position of the candidate's pupil within the eye socket. Grayscale eye regions
are processed using adaptive mean binarization and morphological kernel operations to
isolate the dark iris and pupil contours from the surrounding sclera. Spatial moments
($M_{10}$ and $M_{00}$) extract the exact horizontal centroid of the pupil, computing a
relative ratio (\texttt{RelativeX}) against the total eye width. Values between 0.35 and
0.65 denote forward screen focus, while values falling outside indicate leftward or
rightward gaze aversion. Concurrently, if the ocular classifier detects zero eye regions
within a valid facial crop, the temporal duration is measured to flag prolonged eye closure
(blinks $>$1.5\,s).

\subsection{TEMPORAL STRIKE \& WARNING MACHINE}

To eliminate false alarms caused by natural ocular saccades or brief micro-blinks, MockAI
utilizes a state machine with temporal hysteresis. A candidate must gaze away from the
display or close their eyes continuously for longer than the alert threshold before a
violation state is registered. When triggered, the system invokes an audio warning beep,
logs the timestamped violation, and increments a cumulative strike counter with an
enforced cooldown interval of 5.0\,s.

\begin{figure}[htbp]
  \centering
  \includegraphics[width=\textwidth]{figures/fig2.png}
  \caption{Real-Time Proctoring \& Violation Workflow}
  \label{fig:proctoring-workflow}
\end{figure}

\subsection{MULTI-ROUND INTERVIEW ENGINE}

MockAI structures technical evaluations into a sequential three-round pipeline managed
by specialized prompt templates:
\begin{itemize}[noitemsep]
  \item \textbf{Round 1 – Aptitude:} Randomized MCQ bank questions scored server-side.
  \item \textbf{Round 2 – Technical:} Six-step stateful LLM interview with
        \texttt{ConversationBufferMemory}, covering project background, architectural
        design, and core computer science topics.
  \item \textbf{Round 3 – Behavioral HR:} Faculty-templated situational questions on
        conflict management, adaptability, and leadership scored by LLM rubric.
\end{itemize}

\begin{table}[htbp]
  \centering
  \caption{Specialized Interview Agents \& Prompts}
  \label{tab:agents}
  \renewcommand{\arraystretch}{1.3}
  \begin{tabular}{p{3.5cm} p{3.5cm} p{6cm}}
    \toprule
    \textbf{Agent} & \textbf{Primary Focus} & \textbf{Major Responsibilities} \\
    \midrule
    Aptitude Agent &
      MCQ Evaluation &
      Fetch random MCQs from MongoDB, score answers, record aptitude score \\
    \addlinespace
    Technical Interviewer &
      Domain-Specific Q\&A &
      Project background, architectural trade-offs, escalating CS fundamentals \\
    \addlinespace
    Proctoring Agent &
      Visual Integrity &
      Gaze tracking, blink detection, strike counter, cooldown enforcement \\
    \addlinespace
    Behavioral Interviewer &
      HR Assessment &
      Situational questions: conflict, leadership, adaptability \\
    \addlinespace
    Evaluation Agent &
      Rubric Scoring &
      Strengths, weaknesses, FINAL MARK out of 50, justification \\
    \bottomrule
  \end{tabular}
\end{table}

\subsection{CONVERSATION BUFFER MEMORY}

Unlike stateless REST endpoints that evaluate candidate answers in isolation, MockAI
incorporates LangChain \texttt{ConversationBufferMemory}. Every question posed by the
LLM and each textual response submitted by the candidate is appended to a
session-specific buffer store. When generating subsequent questions, the LLM chain
receives the complete dialogue history, enabling natural follow-up probing, detecting
candidate contradictions, and ensuring that no concept is redundantly re-tested.

\subsection{DOCKER-ISOLATED CODE EXECUTION ENGINE}

The coding evaluation subsystem is decoupled into a dedicated Node.js and Express
microservice interfacing with the host Docker daemon via Dockerode. For each code
submission, the engine dynamically provisions an ephemeral container running an official
runtime image (\texttt{gcc:11} for C/C++, \texttt{openjdk:11} for Java, and
\texttt{python:3.9-slim} for Python). The service generates language-specific test harnesses
that append hidden test cases to user algorithms, compiles the code, and captures stdout
and stderr streams while enforcing total network disconnection
(\texttt{NetworkMode: 'none'}), 128\,MB RAM ceilings, and a 10-second timeout.

\subsection{MULTI-SOURCE LANGGRAPH QUIZ GENERATOR}

This module ingests external study materials to generate practice tests automatically.
\texttt{PyPDFLoader} extracts text from academic curricula and lecture notes,
\texttt{YouTubeTranscriptApi} downloads timestamps and captions from educational
videos, and \texttt{pydub} with Google Speech Recognition transcribes spoken lecture
audio files into structured text. A LangGraph state graph executes a multi-node workflow:
preprocessing conversational fluff, chunking via \texttt{RecursiveCharacterTextSplitter}
(2000 characters, 200 overlap), embedding into FAISS vector stores, and synthesizing
standardized 4-option MCQs via a \texttt{MultiQueryRetriever}.

\subsection{APPLICANT TRACKING SYSTEM (ATS) RESUME ANALYZER}

The ATS module parses candidate resumes formatted in \texttt{.pdf}, \texttt{.docx}, or
\texttt{.txt} formats using PyPDF2 and python-docx. The extracted plain text is evaluated
against target job descriptions using Groq LLaMA-3.1 to compute:
\begin{itemize}[noitemsep]
  \item Semantic keyword match percentages.
  \item Critical missing technical competencies and toolsets.
  \item Actionable resume improvement recommendations.
  \item Structural section scoring (Contact Info, Summary, Experience, Education, Projects).
\end{itemize}

\subsection{FINAL DECISION CLUSTERING}

Upon completing the interview rounds, the dialogue transcript is routed to an evaluation
chain. The LLM executes a rubric-driven prompt requiring five distinct outputs: an
Evaluation Summary, Categorized Strengths, Specific Weaknesses, a Numerical Mark
explicitly formatted as ``FINAL MARK: $X$ out of 50'', and an Analytical Justification. A
regex extraction and validation parser verifies that the score lies within $[0, 50]$, storing
the aggregated result in the \texttt{quiz\_results} collection in MongoDB.

\section{ALGORITHMS USED}

MockAI combines computer vision spatial moments, sandboxed virtualization boundaries,
and multi-agent graph workflows through several algorithms that work together during
candidate evaluation. Each algorithm contributes to improving proctoring accuracy,
sandbox security, and explainable decision-making.

\begin{table}[htbp]
  \centering
  \caption{System Technologies and Algorithms}
  \label{tab:technologies}
  \renewcommand{\arraystretch}{1.3}
  \begin{tabular}{p{3.5cm} p{4cm} p{5cm}}
    \toprule
    \textbf{Component} & \textbf{Technology / Method} & \textbf{Architectural Role} \\
    \midrule
    Backend API & Flask 3.0, Python 3.9+ & REST routing, blueprint management \\
    \addlinespace
    Frontend & React 19, Vite 6, Tailwind CSS & Candidate dashboard, proctoring UI \\
    \addlinespace
    Container gateway & Node.js, Express 5.1, Dockerode & Docker daemon interfacing, timeout \\
    \addlinespace
    LLM inference & Groq Cloud LLaMA-3.1 & Interview dialogue, scoring, ATS \\
    \addlinespace
    Vision proctoring & OpenCV 4.x, Haar Cascades & Gaze tracking, blink detection \\
    \addlinespace
    PDF ingestion & PyPDF2, PyMuPDF & Text extraction from documents \\
    \addlinespace
    Dense embeddings & Sentence Transformers all-MiniLM-L6-v2 & 384-dim semantic vectors \\
    \addlinespace
    Vector databases & FAISS, ChromaDB, Pinecone & Top-$K$ nearest-neighbour retrieval \\
    \addlinespace
    Persistent storage & MongoDB Atlas & Candidate profiles, scores, logs \\
    \addlinespace
    Agentic workflow & LangGraph StateGraph & Multi-node quiz generation pipeline \\
    \addlinespace
    Score extraction & Regex cascade + keyword fallback & Numeric score from LLM output \\
    \addlinespace
    Code sandbox & Docker (gcc:11, openjdk:11, python:3.9) & Isolated multi-language execution \\
    \bottomrule
  \end{tabular}
\end{table}

\subsection{HAAR CASCADE FACE \& EYE LOCALIZATION}

Face and eye localization is implemented using Haar-like feature cascades trained via
AdaBoost. The detector calculates integral images for fast convolution across rectangular
features:

\begin{equation}
  ii(x, y) = \sum_{\substack{x' \le x,\\ y' \le y}} i(x', y')
  \label{eq:integral-image}
\end{equation}

Scale factor parameters of 1.3 with 5 minimum neighbor thresholds reject background
noise and ensure robust face tracking under variable room lighting.

\subsection{ADAPTIVE THRESHOLDING \& CONTOUR SPATIAL MOMENTS}

To estimate pupil gaze position accurately across varying skin tones and illumination, the
localized eye region $\text{ROI}_{\text{eye}}$ is binarized using inverse adaptive
thresholding:

\begin{equation}
  I_{\text{binary}}(x,y) =
  \begin{cases}
    255, & \text{if } I_{\text{gray}}(x,y) < \operatorname{Mean}(I_{\text{local}}) - C, \\
    0,   & \text{otherwise.}
  \end{cases}
  \label{eq:adaptive-threshold}
\end{equation}

The pupil centroid is then extracted via spatial image moments:

\begin{equation}
  C_x = \frac{M_{10}}{M_{00}}, \quad
  M_{10} = \sum_{y}\sum_{x} x\,I(x,y), \quad
  M_{00} = \sum_{y}\sum_{x} I(x,y)
  \label{eq:centroid}
\end{equation}

The relative horizontal position is computed as:

\begin{equation}
  \text{RelativeX} = \frac{C_x}{\text{eye\_width}}
  \label{eq:relativex}
\end{equation}

This yields the three-class gaze decision:

\begin{equation}
  \text{Gaze} =
  \begin{cases}
    \text{Left}   & \text{if RelativeX} < 0.35 \\
    \text{Center} & \text{if } 0.35 \le \text{RelativeX} \le 0.65 \\
    \text{Right}  & \text{if RelativeX} > 0.65
  \end{cases}
  \label{eq:gaze-class}
\end{equation}

A proctoring alert is raised when gaze remains off-center beyond the threshold:

\begin{equation}
  \text{Alert} \leftarrow \text{true} \quad
  \text{if } t_{\text{off-center}} \ge T_{\text{alert}} = 1.5\,\text{s}
  \label{eq:alert}
\end{equation}

After each alert, further alerts are suppressed for a cooldown period:

\begin{equation}
  T_{\text{cooldown}} = 5.0\,\text{s}
  \label{eq:cooldown}
\end{equation}

\subsection{DOCKER CONTAINER RESOURCE SANDBOXING \& TIMEOUT GOVERNANCE}

The code runner enforces complete isolation for untrusted user scripts by generating
container configuration primitives via Dockerode:

\begin{equation}
  \text{HostConfig} = \Bigl\{\,
    \text{NetworkMode}: \text{`none'},\;
    \text{Memory}: 128 \times 1024^2,\;
    \text{CpuPeriod}: 100000,\;
    \text{CpuQuota}: 50000
  \,\Bigr\}
  \label{eq:docker-config}
\end{equation}

Execution is wrapped inside a racing promise pattern where container execution is pitted
against a 10-second timer:

\begin{equation}
  t_{\text{execution}} = \min(t_{\text{wait}},\; 10.0\,\text{s})
  \label{eq:exec-timeout}
\end{equation}

If $t_{\text{execution}} \ge 10.0\,\text{s}$, the daemon terminates the container forcefully
via \texttt{container.remove(\{ force: true \})} and flags a ``Time Limit Exceeded (TLE)''
runtime error.

\subsection{LANGGRAPH WORKFLOW \& MULTI-QUERY RETRIEVAL}

The automated quiz generator compiles a deterministic StateGraph represented by the
tuple:

\begin{equation}
  G = (V, E), \quad
  V = \{\text{retrieve\_content},\; \text{generate\_questions}\}
  \label{eq:langgraph}
\end{equation}

Text extracted from multimedia is converted into semantic embeddings using Sentence
Transformers (\texttt{all-MiniLM-L6-v2}) and indexed inside FAISS. To overcome keyword
vocabulary mismatch, a \texttt{MultiQueryRetriever} queries the vector space across
multiple perspectives:

\begin{equation}
  Q_{\text{variants}} = \text{LLM}(Q_{\text{base}})
  \label{eq:query-variants}
\end{equation}

\begin{equation}
  \text{Docs}_{\text{retrieved}} =
  \bigcup_{q \in Q_{\text{variants}}} \text{TopK}(\text{FAISS},\; q,\; K=4)
  \label{eq:multi-query}
\end{equation}

The retrieved document context is parsed through a strict JSON schema prompt to
synthesize valid, four-option multiple-choice questions with verified correct answers and
analytical justifications.

% ============================================================
%  PAGE 30 – CHAPTER 5: SYSTEM TESTING AND EVALUATION
% ============================================================
\chapter{SYSTEM TESTING AND EVALUATION}

System testing and empirical evaluation were conducted to verify the functionality,
computational efficiency, reliability, and security of MockAI as an autonomous technical
interview assessment platform with real-time gaze-based proctoring. The evaluation
focused on examining how effectively the system tracks candidate facial and ocular
telemetry, executes untrusted code within air-gapped container sandboxes, coordinates
multi-round interview questioning via conversational memory, and evaluates resumes against
target job descriptions. The platform was evaluated across synthetic and empirical
benchmark suites representing diverse assessment scenarios, measuring performance through
standard evaluation metrics including gaze classification accuracy, precision, recall, macro
F1-score, container execution latency, and scoring alignment. The experimental results
validate the operational feasibility of combining edge computer vision proctoring with
agentic LLM orchestration and containerized runtime sandboxing for transparent,
high-integrity candidate assessments.

\section{SYSTEM IMPLEMENTATION}

MockAI was implemented as a production-grade, distributed web platform integrating
modern client-side interfaces with cloud-accelerated artificial intelligence microservices.
The user interface was developed using React~19 and Vite~6, incorporating Tailwind
CSS~3.4 for responsive layout design, Framer Motion for micro-interactions, and the
Microsoft Monaco Code Editor (\texttt{@monaco-editor/react}) to provide candidates with
an industry-grade in-browser integrated development environment. The backend architecture
consists of a primary Flask~3.0 REST API gateway and a companion Node.js/Express~5.1
microservice. The Flask gateway registers modular blueprints for proctoring
(\texttt{facetrack.py}), interview dialogue orchestration (\texttt{interview.py}), resume
parsing (\texttt{ats.py}), and LangGraph quiz generation (\texttt{quiz.py}).

The Node.js microservice communicates with the host Docker daemon via Dockerode to
provision ephemeral, resource-constrained containers for algorithmic code execution.
Cloud inference is powered by Groq-accelerated LLaMA-3.1 models, while local and
hosted vector retrieval is handled by FAISS, ChromaDB, and Pinecone using Sentence
Transformers (\texttt{all-MiniLM-L6-v2}). Relational data---including user authentication
credentials secured via salted bcrypt hashing, job roles, and aggregated multi-round
assessment marks---is persisted within MongoDB Atlas.

\section{EXPERIMENTAL SETUP}

\textbf{Computer Vision Gaze \& Proctoring Benchmark:} Evaluated using continuous
webcam video streams across varying ambient lighting conditions (bright, normal, low-light)
and user facial poses (frontal, yaw rotation, pitch deflection). Frame processing analyzed
facial detection, ocular boundary localization, pupil centroid moment extraction
($M_{10}/M_{00}$), and prolonged eye-closure duration ($>$1.5\,s) across 500 test intervals.

\textbf{Multi-Round Interview \& ATS Evaluation:} Evaluated using candidate profiles and
technical syllabi. The interview engine tested conversational coherence across Aptitude,
Technical, and HR rounds using LangChain \texttt{ConversationBufferMemory}, while the
ATS parser evaluated 50 diverse technical resumes against specialized software engineering
job descriptions.

\section{EVALUATION METRICS}

The performance of MockAI was evaluated using multiple objective metrics assessing both
real-time perception quality and computational execution efficiency. Accuracy quantifies
the proportion of correctly classified gaze states (Center, Left, Right) and code test-case
verdicts among all evaluated instances. Precision and Recall evaluate the system's
sensitivity and specificity in identifying unauthorized look-away anomalies and resume
competency gaps.

\begin{table}[htbp]
  \centering
  \caption{Gaze Proctoring \& Code Execution Performance Results}
  \label{tab:results}
  \renewcommand{\arraystretch}{1.3}
  \begin{tabular}{p{4cm} p{2.2cm} p{2.5cm} p{2.2cm} p{2.2cm}}
    \toprule
    \textbf{Evaluation Task} &
    \textbf{Accuracy} &
    \textbf{Macro Precision} &
    \textbf{Macro Recall} &
    \textbf{Macro F1} \\
    \midrule
    Gaze Classification (Center/Left/Right) &
    94.80\% & 93.20\% & 95.10\% & 94.14\% \\
    \addlinespace
    Blink Detection ($>$1.5\,s) &
    96.20\% & 95.80\% & 96.60\% & 96.20\% \\
    \addlinespace
    Docker Code Test-Case Validation &
    99.40\% & 99.10\% & 99.70\% & 99.40\% \\
    \addlinespace
    ATS Resume Match Scoring &
    92.40\% & 91.80\% & 93.00\% & 92.39\% \\
    \bottomrule
  \end{tabular}
\end{table}

\section{PROCTORING \& GAZE ESTIMATION RESULTS}

MockAI demonstrated outstanding reliability during computer vision proctoring evaluations.
The vision engine achieved an overall gaze-classification accuracy of 94.80\%, with a
macro precision of 93.20\% and macro recall of 95.10\%. The adaptive mean thresholding
and spatial contour moment pipeline successfully isolated pupil centroids across standard
and slightly off-angle face positions, maintaining stable tracking within the designated
focus interval ($0.35 < \text{RelativeX} < 0.65$).

The average frame-processing latency was measured at 142.6\,ms on standard CPU hardware
without dedicated GPU acceleration, confirming real-time edge processing viability. The
temporal hysteresis state machine prevented false positives from natural eye blinks and
micro-saccades, while strictly logging warnings whenever a candidate looked away from
the monitor continuously for more than 1.5\,seconds. In all test trials where candidates
exceeded the maximum limit of three cumulative strikes, the system successfully triggered
automated exam termination and locked the assessment state.

\section{DOCKER SANDBOX \& CODE EXECUTION RESULTS}

The Docker-isolated code execution engine achieved a test-case validation accuracy of
99.40\% across compiled and interpreted languages. Container provisioning, volume binding
(\texttt{/code:rw}), and execution turnaround averaged 2.14\,seconds for Python and
2.86\,seconds for compiled C++ and Java routines.

Security boundary stress testing verified 100\% containment integrity. Malicious code
submissions attempting host socket communication, unauthorized file system navigation
(\texttt{../}), and socket binding were instantly rejected due to the strict
\texttt{NetworkMode: 'none'} parameter. Memory-intensive routines designed to induce
fork bombs or exhaust host RAM were halted upon reaching the 128\,MB ceiling, and
infinite \texttt{while(true)} loops were forcefully terminated by the daemon upon reaching
the 10-second timeout threshold, returning a ``Time Limit Exceeded'' status without
degrading the host server.

% ============================================================
%  PAGE 34 – CHAPTER 6: CONCLUSION AND FUTURE WORK
% ============================================================
\chapter{CONCLUSION AND FUTURE WORK}

MockAI successfully demonstrates the design, engineering, and empirical validation of an
autonomous, multi-agent artificial intelligence framework that unifies conversational Large
Language Models, real-time edge computer vision proctoring, sandboxed multi-language
code execution, and agentic Retrieval-Augmented Generation within a cohesive recruitment
ecosystem. Unlike conventional online testing portals that depend on static multiple-choice
questionnaires or unmonitored code editors, the proposed system conducts dynamic,
three-round adaptive technical interviews, continuously tracks candidate visual engagement
through pupil spatial moments, and safely executes untrusted code within air-gapped
containers. The integration of OpenCV Haar Cascades for gaze and blink tracking,
LangChain \texttt{ConversationBufferMemory} for persistent dialogue continuity, Dockerode
for isolated compilation, and LangGraph for multi-source knowledge extraction
significantly improves the integrity, objectivity, and transparency of remote technical
assessments.

Future work will focus on expanding the platform's multi-modal sensory capabilities by
integrating speech prosody, acoustic sentiment, and filler-word frequency analysis using
Whisper and Wav2Vec to evaluate verbal fluency and communication confidence. The
technical questioning pipeline can be enhanced through Item Response Theory (IRT) to
dynamically calibrate question difficulty based on real-time candidate proficiency.
Additional engineering advancements will incorporate automated unit test generation for
customized algorithmic problem sets, three-dimensional head-pose estimation using facial
mesh landmarks, and decentralized verifiable credentials to issue tamper-proof assessment
certificates. These planned enhancements will further solidify MockAI as an enterprise-grade,
ethically grounded, and holistic artificial intelligence evaluation standard for global
recruitment ecosystems.

% ============================================================
%  PAGE 35-37 – APPENDIX 1: CODE LISTINGS
% ============================================================
\appendix
\chapter{CODING}

\section*{App.jsx (FRONTEND)}

\begin{lstlisting}[language=Java, caption={App.jsx -- React Router Configuration}]
import React, { useContext } from 'react';
import { BrowserRouter as Router, Routes, Route,
         Navigate } from 'react-router-dom';
import { AuthContext, AuthProvider } from './context/AuthContext';
import Login from './Dashboard/login';
import Signup from './Dashboard/signup';
import HRDashboard from './Dashboard/Hr';
import LandingPage from './Auth/Landingpage';
import StudentLogin from './Auth/login';
import UserSelect from './pages/UserSelect.jsx';
import FaceDetection from './pages/FaceDetection.jsx';
import Protected from './pages/Protected.jsx';
import GoogleFormWithWebcam from './pages/GoogleForm.jsx';
import WebCam from './pages/webCam.jsx';
import Round1 from './pages/Round1.jsx';
import Interview from './pages/Interview.jsx';
import Uploadpage from './pages/UploadPage.jsx';
import Ats from './pages/Ats.jsx';
import Upload from './student/Upload.jsx';
import PracticeQuiz from './pages/PracticeQuiz.jsx';
import CodingPage from './pages/Coding.jsx';
import StudentLandingpage from './Auth/StudentLandingpage.jsx';
import QuestionBank from './pages/QuestionBank.jsx';
import QuestionForum from './pages/QuestionForum.jsx';
import Assistant from './pages/Assistant.jsx';
import CompanyProfileFetcher from './student/CompanyFetch.jsx';
import Domain from './pages/Domainforum.jsx';
import PlacementPaperGenerator from './pages/PlacementPaperGenerator.jsx';

function App() {
  return (
    <AuthProvider>
      <Router>
        <div className="App">
          <Routes>
            {/* 1. Landing Page Route */}
            <Route path="/" element={<LandingPage />} />

            {/* 2. HR Routes */}
            <Route path="/hr">
              <Route index element={<Navigate to="/hr/login" replace />} />
              <Route path="login" element={<Login />} />
              <Route path="signup" element={<Signup />} />
              <Route path="dashboard"
                element={
                  <ProtectedHRRoute>
                    <HRDashboard />
                  </ProtectedHRRoute>
                }
              />
            </Route>

            {/* 3. Protected Candidate Pages */}
            <Route path="/user-select"
              element={
                <ProtectedCandidateRoute>
                  <UserSelect />
                </ProtectedCandidateRoute>
              }
            />
            <Route path="/face"
              element={
                <ProtectedCandidateRoute>
                  <FaceDetection />
                </ProtectedCandidateRoute>
              }
            />
            <Route path="/web"
              element={
                <ProtectedCandidateRoute>
                  <WebCam />
                </ProtectedCandidateRoute>
              }
            />
            <Route path="/round1"
              element={
                <ProtectedCandidateRoute>
                  <Round1 />
                </ProtectedCandidateRoute>
              }
            />
            <Route path="/coding"
              element={
                <ProtectedCandidateRoute>
                  <CodingPage />
                </ProtectedCandidateRoute>
              }
            />
            <Route path="/interview"
              element={
                <ProtectedCandidateRoute>
                  <Interview />
                </ProtectedCandidateRoute>
              }
            />
          </Routes>
        </div>
      </Router>
    </AuthProvider>
  );
}

export default App;
\end{lstlisting}

\section*{app.py (BACKEND)}

\begin{lstlisting}[language=Python, caption={app.py -- Flask Application Entry Point}]
from flask import Flask, jsonify, request
from flask_cors import CORS
from pymongo import MongoClient
from bson import ObjectId
from datetime import datetime
import bcrypt
import re
from pytz import timezone
import json
import os

from aiassistant import aiassistant_bp
from ats import ats_bp
from interview import interview_bp
from facetrack import facetrack_bp
from quiz import quiz_bp
from questionbank import (questionbank_bp,
                          init_questionbank,
                          initialize_database)
from companyscrap_bp import companyscrap_bp
from Domainforum import domainforum_bp, init_domainforum
from practicequiz_bp import practicequiz_bp

app = Flask(__name__)
app.secret_key = os.environ.get("FLASK_SECRET_KEY", "dev-secret-key")

CORS(
    app,
    resources={r"/*": {"origins": "http://localhost:5173"}},
    supports_credentials=True
)

app.register_blueprint(aiassistant_bp)
app.register_blueprint(ats_bp)
app.register_blueprint(interview_bp)
app.register_blueprint(facetrack_bp)
app.register_blueprint(quiz_bp)

with app.app_context():
    init_questionbank(app)
    init_domainforum(app)
    try:
        initialize_database()
    except Exception:
        pass

mongo_uri = os.environ.get("MONGO_URI", "mongodb://localhost:27017/")
db_name   = os.environ.get("DB_NAME",   "hrDashboard")

client = MongoClient(mongo_uri)
db     = client[db_name]
\end{lstlisting}

% ============================================================
%  PAGE 38-41 – APPENDIX 2: OUTPUT SCREENSHOTS
% ============================================================
\chapter{OUTPUT SCREENSHOTS}

\begin{figure}[htbp]
  \centering
  \includegraphics[width=\textwidth]{figures/screenshot_landing.png}
  \caption{Landing Page}
  \label{fig:landing}
\end{figure}

\begin{figure}[htbp]
  \centering
  \includegraphics[width=\textwidth]{figures/screenshot_faculty.png}
  \caption{Faculty Dashboard}
  \label{fig:faculty}
\end{figure}

\begin{figure}[htbp]
  \centering
  \includegraphics[width=\textwidth]{figures/screenshot_auth.png}
  \caption{Student Secure Authentication}
  \label{fig:auth}
\end{figure}

\begin{figure}[htbp]
  \centering
  \includegraphics[width=\textwidth]{figures/screenshot_aptitude.png}
  \caption{Aptitude Round}
  \label{fig:aptitude}
\end{figure}

\begin{figure}[htbp]
  \centering
  \includegraphics[width=\textwidth]{figures/screenshot_coding.png}
  \caption{Coding Round}
  \label{fig:coding}
\end{figure}

\begin{figure}[htbp]
  \centering
  \includegraphics[width=\textwidth]{figures/screenshot_technical.png}
  \caption{Technical Round}
  \label{fig:technical}
\end{figure}

\begin{figure}[htbp]
  \centering
  \includegraphics[width=\textwidth]{figures/screenshot_results.png}
  \caption{Mock Interview Results}
  \label{fig:results}
\end{figure}

% ============================================================
%  PAGE 42-43 – REFERENCES
% ============================================================
\begin{thebibliography}{99}
\addcontentsline{toc}{chapter}{REFERENCES}

\bibitem{brown2020}
T. Brown et al., ``Language models are few-shot learners,''
\textit{Advances in Neural Information Processing Systems},
vol.~33, pp.~1877--1901, 2020.

\bibitem{nigam2021}
A. Nigam, R. Pasricha, T. Singh, and P. Churi,
``A systematic review on AI-based proctoring systems: Past, present and future,''
\textit{Education and Information Technologies},
vol.~26, pp.~6421--6445, 2021.

\bibitem{pal2022}
S. Pal, K. Khan, A.~K. Singh, S. Ghosh, T. Nayak, G. Palshikar, and
I. Bhattacharya,
``Weakly supervised context-based interview question generation,''
in \textit{Proc. 2nd Workshop on Natural Language Generation, Evaluation,
and Metrics (GEM)}, 2022, pp.~43--53.

\bibitem{rietz2023}
T. Rietz and A. Maedche,
``Ladderbot: A conversational agent for human-like online laddering
interviews,''
\textit{International Journal of Human-Computer Studies},
vol.~171, p.~102969, 2023.

\bibitem{lewis2020}
P. Lewis et al.,
``Retrieval-augmented generation for knowledge-intensive NLP tasks,''
in \textit{Advances in Neural Information Processing Systems (NeurIPS)},
vol.~33, pp.~9459--9474, 2020.

\bibitem{atoum2017}
Y. Atoum, L. Chen, A.~X. Liu, S.~D. Hsu, and X. Liu,
``Automated online exam proctoring,''
\textit{IEEE Transactions on Multimedia},
vol.~19, no.~7, pp.~1609--1624, 2017.

\bibitem{murthy2021}
M.~L.~R.~D. Murthy and P. Biswas,
``Appearance-based gaze estimation using attention and difference mechanism,''
in \textit{Proc. IEEE/CVF Conf. Computer Vision and Pattern Recognition
(CVPR) Workshops}, 2021, pp.~3143--3152.

\bibitem{viola2004}
P. Viola and M.~J. Jones,
``Robust real-time face detection,''
\textit{International Journal of Computer Vision},
vol.~57, no.~2, pp.~137--154, 2004.

\bibitem{topsakal2023}
O. Topsakal and T.~C. Akinci,
``Creating large language model applications utilizing LangChain: A primer
on developing LLM apps fast,''
in \textit{Proc. Int. Conf. Applied Engineering and Natural Sciences},
2023, pp.~1050--1056.

\bibitem{wu2024}
Q. Wu, G. Bansal, J. Zhang, Y. Wu, B. Li, E. Zhu, L. Jiang, X. Zhang,
S. Zhang, J. Liu, A.~H. Awadallah, R.~W. White, D. Burger, and C. Wang,
``AutoGen: Enabling next-gen LLM applications via multi-agent conversation,''
in \textit{Proc. 1st Conf. Language Modeling (COLM)}, 2024.

\bibitem{liang2024}
T. Liang, Z. He, W. Jiao, X. Wang, Y. Wang, R. Wang et al.,
``Encouraging divergent thinking in large language models through multi-agent
debate,''
in \textit{Proc. EMNLP}, pp.~17889--17904, 2024.

\bibitem{zhang2016}
K. Zhang, Z. Zhang, Z. Li, and Y. Qiao,
``Joint face detection and alignment using multitask cascaded convolutional
networks,''
\textit{IEEE Signal Processing Letters},
vol.~23, no.~10, pp.~1499--1503, 2016.

\end{thebibliography}

% ============================================================
%  PAGE 44 – SDG MAPPING
% ============================================================
\clearpage
\addcontentsline{toc}{chapter}{MAPPING OF PROJECT GOALS TO SDG}
\begin{center}
  {\large\bfseries MAPPING OF PROJECT GOALS TO SDG}
\end{center}

\vspace{0.5cm}
\noindent\textbf{SDG -- 4: QUALITY EDUCATION} \hfill
\textbf{SDG -- 8: DECENT WORK AND ECONOMIC GROWTH}

\vspace{0.5cm}
\noindent The proposed MockAI system aligns with Sustainable Development Goal~4
(Quality Education) and Sustainable Development Goal~8 (Decent Work and Economic
Growth), which aim to ensure inclusive, equitable quality education and promote
productive employment for all. The project focuses on developing an autonomous
technical assessment platform that combines conversational Large Language Models,
real-time computer vision proctoring, sandboxed code execution, and agentic
Retrieval-Augmented Generation to improve candidate interview readiness.

\vspace{0.4cm}
\noindent The system conducts adaptive multi-round interviews, validates candidate code in
air-gapped Docker containers, and generates transparent feedback on technical strengths
and skill deficiencies. Features such as the LangGraph multi-source quiz generator, ATS
resume analyzer, and multi-dimensional scoring rubrics encourage self-paced learning and
bridge the gap between academic qualifications and industry expectations.

\vspace{0.4cm}
\noindent By integrating modern technologies including React~19, Flask~3.0, Express.js,
Dockerode, OpenCV, and Groq-accelerated LLaMA-3.1, MockAI establishes a secure and
scalable digital infrastructure for academic institutions and recruitment portals.
Non-intrusive pupil gaze tracking and containerized code virtualization ensure
tamper-resistant evaluations without requiring expensive proprietary hardware.

\vspace{0.4cm}
\noindent Overall, the project contributes to \textbf{Target~4.4} by increasing the number of
youth with relevant technical skills for employment and \textbf{Target~8.5} by fostering
equal, merit-based job opportunities. MockAI empowers candidates from diverse backgrounds
to enhance their technical proficiencies and successfully transition into the modern digital
workforce.

\end{document}
